% sections/liquidity.tex: Section 5.6: default as a choice at a margin; the payment margin under local liquidity shocks.
% Results: cross_loan_status.csv (E3), level_shock_annuity.csv (L3), panel_windows_auto.csv (cohorts).
\subsection{Liquidity Shocks and the Payment Margin}\label{sec:liquidity}

Default that is forced by a collapse of cash flow rather than chosen would enter the payment margin without revealing a preference. Three facts establish that the estimates are choices at a margin, and the formula's own treatment of liquidity, through the marginal-utility ratio, completes them. Table \ref{tab:liquidity} in Appendix \ref{appendix:figures} reports the first two.

First, auto default is overwhelmingly a portfolio choice. Among 2.46 million auto defaulters holding other open, clean tradelines at the time of default, 65 percent are still making payments on at least one other account six months after the auto charge-off (they hold 10.7 other tradelines on average), and only 1.6 percent default on every account. Money continued to flow to other creditors through the event. The only form of the objection that undermines revealed preference, a total collapse of cash flow, is bounded between 1.6 and roughly 35 percent of defaults.

Second, the payment margin does not track liquidity shocks. Sorting loans within origination year into terciles of the change in state unemployment over the loan's first eighteen months, the logit coefficient of default within twelve months on the log payment is $0.53$, $0.62$ and $0.61$ from the mildest to the most severe exposure (Table \ref{tab:liquidity}, Panel B), and $0.57$ against $0.52$ for expansion and recession cohorts. On the population specification the elasticity of charge-off within 24 months is $0.40$, $0.40$, $0.38$ and $0.37$ for loans originated in 2013--15, 2016--18, 2019--21 and 2022 and after, cohorts that met the pandemic and the 2022 rate cycle at different loan ages, and $0.39$ for 2010--12 (Table \ref{tab:exposure} in Appendix \ref{appendix:figures}). If the margin repackaged the shocks a cohort met, it would move with them.

Third, distress is persistent but not absorbing, and liquidity enters the formula where the theory puts it. One fifth of delinquent auto trades cure within a four-month archive gap (Appendix \ref{appendix:measuredrho}), so the marginal defaulter is in a high-marginal-utility state that she expects to leave; that is the premium $\lambda$, which Section \ref{sec:quasi} measures at $0.2$ rather than assuming away. A liquidity constraint that binds in a given month raises the density of borrowers at the default threshold, the level of the payment elasticity in which the heterogeneity of Section \ref{sec:heterogeneity} appears; the rate is a property of the weights, which the level does not enter.
